\begin{abstract}
Diffusion models have emerged as the dominant paradigm for high-fidelity image generation, yet their training remains
computationally expensive and slow to converge. Recent work introduced REPA (REPresentation Alignment), a regularization
technique that aligns projections of noisy input hidden states in diffusion transformers with clean image representations
from pretrained visual encoders. Reported as an ICLR 2025 Oral, REPA demonstrated substantial improvements in training
efficiency, achieving up to 17.5x faster convergence on large-scale ImageNet experiments, while also
improving final generation quality.
In this project, we aim to present a partial reproduction and systematic extension of REPA under constrained computational settings.
Rather than scaling to full ImageNet, we plan to reproduce the core REPA framework on reduced datasets and smaller model configurations, verifying
its key claim of accelerated convergence relative to standard diffusion training. Building on this, we plan to investigate REPA in low-data
regimes, considering both subsampled large-scale datasets (e.g., ImageNet) and naturally small, domain-specific datasets such as Stanford Cars.
We will study three main questions: (1) whether REPA's improvements in convergence speed and sample quality persist when training data is
limited; (2) how the choice of pretrained feature extractor -- comparing alternatives such as DINOv2 and CLIP -- affects performance,
particularly in narrow visual domains; and (3) how sensitive REPA is to dataset scale, including whether its alignment objective can
partially compensate for reduced data availability. Through these experiments, we aim to better characterize the practical benefits
and limitations of representation alignment for efficient diffusion training in realistic, resource-constrained settings.
\end{abstract}
